\chapter{EML Registry: Tamari Lattice Contraction}

\section{Core Types}

\begin{definition}[EML Tree]
\label{def:emltree}
\lean{EMLRegistry.EMLTree}
\leanok
The type of binary trees representing configurations in the Tamari lattice.
A tree is either a \texttt{Leaf} (empty configuration) or a \texttt{Node} $t_1\ t_2$ (composition of two sub-configurations).
\end{definition}

\begin{definition}[Tree Size]
\label{def:size}
\lean{EMLRegistry.EMLTree.size}
\leanok
\uses{def:emltree}
The number of internal nodes in a tree, used as the lattice index $n$ in $T_n$.
\end{definition}

\begin{definition}[Single-Step Contraction]
\label{def:contracts_one}
\lean{EMLRegistry.contracts_one}
\leanok
\uses{def:emltree}
The primitive non-associative rewrite rule: $(a \bullet b) \bullet c \to a \bullet (b \bullet c)$.
This is the coupling signature for choice resolution, lifted through subtrees.
\end{definition}

\begin{definition}[Multi-Step Contraction]
\label{def:contracts_to}
\lean{EMLRegistry.contracts_to}
\leanok
\uses{def:contracts_one}
The reflexive-transitive closure of \textsf{contracts\_one}. Represents an audit trail of choice resolutions with monotonic provenance.
\end{definition}

\begin{definition}[Right-Comb Normal Form]
\label{def:rightComb}
\lean{EMLRegistry.rightComb}
\leanok
\uses{def:emltree}
The minimum element of the Tamari lattice — the equilibrium attractor.
Defined as $\mathsf{rightComb}(0) = \mathsf{Leaf}$ and $\mathsf{rightComb}(n+1) = \mathsf{Node}(\mathsf{Leaf},\ \mathsf{rightComb}(n))$.
\end{definition}

\section{Lifting Lemmas}

\begin{lemma}[Contraction Preserved Under Left Node]
\label{lem:node_left}
\lean{EMLRegistry.contracts_to_node_left}
\leanok
\uses{def:contracts_one, def:contracts_to}
If $t \to t'$ then $\mathsf{Node}(t,\ r) \to \mathsf{Node}(t',\ r)$.
Choice resolution in the left subtree propagates to the whole composition.
\end{lemma}

\begin{lemma}[Contraction Preserved Under Right Node]
\label{lem:node_right}
\lean{EMLRegistry.contracts_to_node_right}
\leanok
\uses{def:contracts_one, def:contracts_to}
If $r \to r'$ then $\mathsf{Node}(l,\ r) \to \mathsf{Node}(l,\ r')$.
Choice resolution in the right subtree propagates to the whole composition.
\end{lemma}

\begin{lemma}[Transitivity of Contraction]
\label{lem:trans}
\lean{EMLRegistry.contracts_to_trans}
\leanok
\uses{def:contracts_to}
If $s \to t$ and $t \to u$ then $s \to u$.
Audit trails concatenate monotonically.
\end{lemma}

\section{Size Invariants}

\begin{lemma}[Size is Preserved by Single-Step Contraction]
\label{lem:size_one}
\lean{EMLRegistry.contracts_one_size_eq}
\leanok
\uses{def:size, def:contracts_one}
Rotation rewrites the tree depth but does not change the number of internal nodes.
\end{lemma}

\begin{lemma}[Size is Preserved by Multi-Step Contraction]
\label{lem:size_multi}
\lean{EMLRegistry.contracts_to_size_eq}
\leanok
\uses{def:size, def:contracts_to, lem:size_one}
Contraction paths preserve tree size, establishing that the Tamari lattice $T_n$ is graded by node count $n$.
\end{lemma}

\section{Main Contraction Theorem}

\begin{lemma}[Node of Right-Combs Contracts to Right-Comb]
\label{lem:node_rightcombs}
\lean{EMLRegistry.node_of_rightCombs_contracts_to_rightComb}
\leanok
\uses{def:rightComb, def:contracts_one, def:contracts_to, lem:node_right}
For any $a, b : \mathbb{N}$,
\[
\mathsf{Node}(\mathsf{rightComb}(a),\ \mathsf{rightComb}(b)) \longrightarrow \mathsf{rightComb}(1 + a + b).
\]
The equilibrium attractor is closed under non-associative composition.
\end{lemma}

\begin{theorem}[Every Tree Contracts to Right-Comb]
\label{thm:contracts_to_rightComb}
\lean{EMLRegistry.contracts_to_rightComb}
\leanok
\uses{def:emltree, def:size, def:rightComb, def:contracts_to, lem:node_left, lem:node_right, lem:trans, lem:node_rightcombs}
For any tree $t : \mathsf{EMLTree}$,
\[
t \longrightarrow \mathsf{rightComb}(t.\mathsf{size}).
\]
Every configuration has a valid evolution path to the equilibrium attractor indexed by its size.
\end{theorem}

\section{Router and Registry}

\begin{definition}[Router Index]
\label{def:router}
\lean{EMLRegistry.RouterIndex}
\leanok
A bounded natural number $\mathsf{Fin}(n)$ serving as a neural binding address.
\end{definition}

\begin{definition}[Type Registry]
\label{def:registry}
\lean{EMLRegistry.TypeRegistry}
\leanok
\uses{def:router, def:emltree}
An injective mapping from router indices to EML trees. This is the cortex-registry interface where neural computation meets formal grammar.
\end{definition}

\section{Witness-Skeptic Game}

\begin{definition}[Cortex Certificate]
\label{def:certificate}
\lean{EMLRegistry.CortexCertificate}
\leanok
\uses{def:contracts_to, def:rightComb}
A proof-carrying audit trail $\langle source,\ target,\ proof \rangle$ showing that a tree reaches its equilibrium.
\end{definition}

\begin{definition}[Certify]
\label{def:certify}
\lean{EMLRegistry.certify}
\leanok
\uses{def:contracts_to, def:rightComb, def:certificate, thm:contracts_to_rightComb}
Issues a \textsf{CortexCertificate} for any tree $t$, certifying that $t \to \mathsf{rightComb}(t.\mathsf{size})$.
This is the quench witness in the Witness-Skeptic game.
\end{definition}

\chapter{Logic Types: Pluralistic Reasoning Framework}

\section{Logic Type Hierarchy}

\begin{definition}[13 Logic Types]
\label{def:logictype}
\lean{LogicTypes.LogicType}
\leanok
An inductive type enumerating 13 reasoning paradigms.

Each handles specific classes of paradoxes as boundary conditions rather than errors.
\end{definition}

\begin{definition}[Logic Classification]
\label{def:logicclass}
\lean{LogicTypes.LogicClass}
\leanok
\uses{def:logictype}
Classification of logic types as Substructural, Extension, Alternative, Generalization, or Classical.
\end{definition}

\begin{definition}[Logic Tree]
\label{def:logictree}
\lean{LogicTypes.LogicTree}
\leanok
\uses{def:emltree, def:logictype}
A type family assigning EML trees to each logic type. Currently uniform; extensible to logic-specific tree types.
\end{definition}

\begin{definition}[Logic Contraction]
\label{def:logiccontraction}
\lean{LogicTypes.LogicContraction}
\leanok
\uses{def:contracts_to, def:logictype}
A type family of contraction relations indexed by \textsf{LogicType}. For Classical logic this is the Tamari contraction from EMLRegistry.
\end{definition}

\begin{definition}[Logic Translation]
\label{def:translation}
\lean{LogicTypes.LogicTranslation}
\leanok
\uses{def:logictype, def:logiccontraction}
A structure specifying forward and backward maps between logic types with soundness and completeness proofs.
\end{definition}

\begin{definition}[Meta-Contraction]
\label{def:metacontracts}
\lean{LogicTypes.MetaContractsTo}
\leanok
\uses{def:logiccontraction, def:translation}
Cross-logic contraction allowing reasoning that spans multiple logic types via intra-logic contraction, inter-logic translation, and transitivity.
\end{definition}

\section{Cayley-Dickson Connection}

\begin{definition}[CD Step]
\label{def:cdstep}
\lean{LogicTypes.LogicType.cdStep}
\leanok
\uses{def:logictype}
Maps logic types to Cayley-Dickson construction steps.
\end{definition}

\begin{definition}[Split-Octonion Sectors]
\label{def:split}
\lean{LogicTypes.LogicType.isAssociativeSector}
\leanok
\uses{def:logictype}
Divides logic types into associative and non-associative sectors of the split-octonion algebra.
\end{definition}

\section{Normal Forms and Contraction Theorem}

\begin{definition}[Logic Normal Form]
\label{def:logicnormal}
\lean{LogicTypes.LogicNormalForm}
\leanok
\uses{def:rightComb, def:logictype}
The normal form for each logic type. For Classical logic this is \textsf{rightComb}.
\end{definition}

\begin{theorem}[Logic Contraction to Normal Form]
\label{thm:logic_contracts}
\lean{LogicTypes.logic_contracts_to_normal_form}
\uses{def:logicnormal, def:logiccontraction, def:size, thm:contracts_to_rightComb}
For any logic type $lt$ and tree $t$, $\mathsf{LogicContraction}\ lt\ t\ (\mathsf{LogicNormalForm}\ lt\ t.\mathsf{size})$.
Currently proven for Classical logic; other cases are placeholders.
\end{theorem}
\chapter{Liar Paradox: The Hard Test Case}

\section{Canonical Encoding}

\begin{definition}[Liar Tree]
\label{def:liar}
\lean{LiarParadox.liar}
\leanok
\uses{def:emltree}
The canonical Liar sentence: a size-3 symmetric tree
\[
\mathsf{Node}(\mathsf{Node}(\mathsf{Leaf},\ \mathsf{Leaf}),\ \mathsf{Node}(\mathsf{Leaf},\ \mathsf{Leaf}))
\]
encoding the self-referential equation $L = \neg L$ via repeated bracketing.
Each leaf pair is structurally identical, so neither bracketing dominates ---
the paradox is that rotation can pass the ``hot potato'' back and forth indefinitely.
\end{definition}

\begin{theorem}[Liar Contracts to Right-Comb]
\label{thm:liar_contracts}
\lean{LiarParadox.liar_contracts_to_rightComb}
\leanok
\uses{def:emltree, thm:contracts_to_rightComb}
\[
\mathsf{liar} \longrightarrow \mathsf{rightComb}(3)
\]
The Liar reaches equilibrium under Classical Tamari contraction, as a special case of the general theorem.
\end{theorem}

\section{Confluence}

\begin{theorem}[Classical Confluence of the Liar]
\label{thm:liar_confluence_classical}
\lean{LiarParadox.liar_confluence_classical}
\leanok
\uses{def:liar, def:contracts_to, def:rightComb, lem:size_multi, thm:contracts_to_rightComb}
If $\mathsf{liar} \to t$ then $t \to \mathsf{rightComb}(3)$.
All Classical contraction paths from the Liar converge to the unique minimum of $T_3$.
\end{theorem}

\begin{theorem}[Cross-Logic Confluence of the Liar (Placeholder)]
\label{thm:liar_confluence_cross}
\lean{LiarParadox.liar_confluence_cross}
\uses{def:logictype, def:logiccontraction, def:liar}
If $\mathsf{liar} \to_{lt_1} t_1$ and $\mathsf{liar} \to_{lt_2} t_2$, there exists a translation between $t_1$ and $t_2$.
This is the \textbf{open universe} property: different logics produce different normal forms,
connected by translations but not necessarily equal.
\textbf{Status:} Returns $\mathsf{True}$ as a placeholder. The real proof requires
\hyperref[sec:missing]{logic-specific contraction relations}, which are not yet defined.
\end{theorem}

\section{Emergent Distance and Friction}

\begin{definition}[Contraction Cost]
\label{def:contraction_cost}
\lean{LiarParadox.contractionCost}
\uses{def:logictype, def:emltree}
The pentagonator distance $\Phi(lt, t)$ for resolving tree $t$ under logic type $lt$.
\textbf{Status:} Stub returning 0 for all logics. A proper implementation requires
a shortest-path search over the contraction graph for each logic type.
\end{definition}

\begin{definition}[Friction Lagrangian]
\label{def:friction_lagrangian}
\lean{LiarParadox.frictionLagrangian}
\uses{def:contraction_cost}
The action functional $L(lt, t) = \Phi(lt, t)$ evaluated on a single
(logic type, tree) pair. \textbf{Status:} Currently identical to \textsf{contractionCost};
the full infinite-series form with commutator and associator derivatives is not yet formalized.
\end{definition}

\begin{definition}[Optimal Resolution]
\label{def:optimal}
\lean{LiarParadox.optimalResolution}
\uses{def:logictype, def:friction_lagrangian}
The logic type minimizing the Friction Lagrangian for a given tree.
\textbf{Status:} Returns $\mathsf{Classical}$ unconditionally. Requires real cost data.
\end{definition}

\section{Missing Work}
\label{sec:missing}

The following definitions and theorems are stated but not yet implemented.
They represent the \textbf{real friction} in the framework ---
the Liar Paradox is defined and its Classical resolution is proven,
but cross-logic analysis cannot proceed until these pieces are populated.

\subsection{Sorries in \textsf{LogicTypes.lean}}

\begin{itemize}
  \item \textbf{\textsf{LogicContraction}} for Fuzzy, ManyValued, Paraconsistent, Temporal, Deontic,
    Epistemic, Quantum, Intuitionistic, Relevance, Free, Infinitary, Modal ---
    only Classical is defined (as the Tamari contraction from EMLRegistry).
  \item \textbf{\textsf{LogicNormalForm}} for all non-Classical logics ---
    only $\mathsf{rightComb}$ (Classical) is defined.
  \item \textbf{\textsf{logic\_contracts\_to\_normal\_form}} for non-Classical logics ---
    the theorem is stated but only the Classical case is proven.
\end{itemize}

\subsection{Stubs in \textsf{LiarParadox.lean}}

\begin{itemize}
  \item \textbf{\textsf{contractionCost}} --- returns 0 for all logics.
    Needs a bounded shortest-path search over the contraction graph of each logic type,
    analogous to $\mathsf{decidable\_contracts\_to}$ but cost-weighted.
  \item \textbf{\textsf{frictionLagrangian}} --- currently the identity on cost.
    The full form $\sum_k \mu_k [e^{\alpha_k C^{(k)}} - \beta_k \ln(1 + (A^{(k)})^2)]$
    with coupling $\sum_k \nu_k (C^{(k)} A^{(k+1)} - A^{(k)} C^{(k+1)})$ is not yet formalized.
  \item \textbf{\textsf{optimalResolution}} --- returns Classical unconditionally.
  \item \textbf{\textsf{liar\_confluence\_cross}} --- returns $\mathsf{True}$; the real proof
    depends on the missing \textsf{LogicContraction} definitions.
\end{itemize}

\subsection{Blocking Graph}

The dependencies among missing work:

\[
\begin{array}{c}
\text{LogicContraction (12 non-Classical)} \\
\downarrow \\
\text{LogicNormalForm (12 non-Classical)} \\
\downarrow \\
\text{logic\_contracts\_to\_normal\_form} \\
\downarrow \\
\text{LiarParadox.contractionCost} \quad \text{LiarParadox.liar\_confluence\_cross} \\
\downarrow \\
\text{LiarParadox.frictionLagrangian} \quad \text{LiarParadox.partitionFunction}
\end{array}
\]

Every missing piece in the chain below \textsf{LogicContraction} is blocked by it.
The Liar formalization makes this blocking structure explicit: previously the sorries
in \textsf{LogicTypes.lean} were abstract ``to be filled later'' --- now they are concrete
dependencies of the Liar's cross-logic analysis.

\endinput